\documentclass[10pt,a4paper]{amsart}
\usepackage[margin=19mm]{geometry}
\usepackage{amsmath,amssymb,mathtools}
\usepackage[scr=rsfs,cal=euler]{mathalfa}
\usepackage{xfrac}
\usepackage[unicode,hidelinks,bookmarks=false]{hyperref}
\newcommand{\ie}{i.e.\ }
\newcommand{\cf}{cf.\ }
\newcommand{\resp}{respectively\ }
\newcommand{\Subsection}{\subsection}
\providecommand{\nobreakdash}{\nobreak\hbox{-}}
\setlength{\emergencystretch}{2em}
\multlinegap0pt
\usepackage{bm}
\usepackage{bbm}
\usepackage{dsfont}
\usepackage{xspace}
\usepackage{cancel}
\usepackage{mathtools}
\usepackage{amsthm}
\makeatletter
\@ifundefined{assumption}{\newtheorem{assumption}{Assumption}}{}
\@ifundefined{lemma}{\newtheorem{lemma}{Lemma}}{}
\@ifundefined{proposition}{\newtheorem{proposition}{Proposition}}{}
\makeatother
\def\convP{\stackrel{\mbox{$\scriptstyle P$}}{\rightarrow}}
\newcommand \ag{\alpha}
\newcommand \eg{\varepsilon}
\newcommand \sg{\sigma}
\newcommand\bZ{{\bf Z}}
\newcommand\bz{{\bf z}}
\newcommand\cZ{{\mathcal Z}}
\newcommand\lbr{\left \{}
\newcommand\lb{\left [}
\newcommand\lip{\left \langle}
\newcommand\lmo{\left |}
\newcommand\lp{\left (}
\newcommand\mbR{{\mathbb R}}
\newcommand\rbr{\right \}}
\newcommand\rb{\right ]}
\newcommand\rip{\right \rangle}
\newcommand\rmo{\right |}
\newcommand\rp{\right )}
\newcommand\var{\hbox{\rm var}}
\title[Extremal correlation coefficient for functional data]{Extremal correlation coefficient for functional data}
\author{Mihyun Kim, Piotr Kokoszka}
\date{Source: arXiv:2405.17318v3 (CC BY 4.0)}
\begin{document}
\maketitle

\noindent\emph{Source and licence.} Extremal correlation coefficient for functional data. Version arXiv:2405.17318v3 (CC BY 4.0), distributed under the Creative Commons Attribution 4.0 International licence, as recorded in the arXiv licence metadata for this exact version and shown on the abstract page https://arxiv.org/abs/2405.17318v3. Full rights notice: licenses/arxiv-2405.17318-CC-BY-4.0.txt.\par\smallskip

\noindent\textit{Editorial scope.} Counted unit: the Proposition whose source label is \texttt{p:bar} in the article (source0.tex lines 1217--1222, source label \texttt{p:bar}), delivered together with the article's own complete original proof of it (source0.tex lines 1224--1297, one \texttt{proof} environment of 74 lines). No mathematics is written, repaired or re-derived: statement and proof are the article's own text line for line. Required context retained uncounted: p:sglim (lines 512--517); a:L2L2 (lines 551--553); l:alM (lines 1047--1052); l:Bern (lines 1182--1200); eq:gi (lines 1185--1187); eq:b (lines 1189--1191); eq:v (lines 1193--1195). Every definition, display and label of those lines is retained unchanged. Omitted from this package and not redistributed: 1--511, 518--550, 554--1046, 1053--1181, 1188--1188, 1192--1192, 1196--1216, 1223--1223, 1298--1851. Nothing inside a retained excerpt is omitted or altered: every retained body is delivered exactly as the article's lines are written, so no omission receipt of any kind accompanies this package. Citations: the retained excerpts carry no citation command at all, so no bibliography entry of the article is transcribed and no reference is redistributed. Reference adaptation: none was needed and none was made; every reference inside the delivered unit resolves inside it, so no unresolved reference or citation is produced. Printed slips kept literal: no printed word, symbol, hypothesis or number of the counted statement or of its proof was altered. Layout and class: the article's own class is \texttt{biometrika} with packages including amsmath, color, newtxtext, newtxmath, algorithm2e; the wrapper is the lane's pinned \texttt{amsart} editorial wrapper at 10pt on A4, which restates the theorem-like environments the retained text needs and adds the article's own macro definitions that the retained text needs (\texttt{\textbackslash ag}, \texttt{\textbackslash bZ}, \texttt{\textbackslash bz}, \texttt{\textbackslash cZ}, \texttt{\textbackslash convP}, \texttt{\textbackslash eg}, \texttt{\textbackslash lb}, \texttt{\textbackslash lbr}, \texttt{\textbackslash lip}, \texttt{\textbackslash lmo}, \texttt{\textbackslash lp}, \texttt{\textbackslash mbR}, \texttt{\textbackslash rb}, \texttt{\textbackslash rbr}, \texttt{\textbackslash rip}, \texttt{\textbackslash rmo}, \texttt{\textbackslash rp}, \texttt{\textbackslash sg}, \texttt{\textbackslash var}), with extra wrapper packages (bm, bbm, dsfont, xspace, cancel, mathtools). The delivered numbering is the unit's own. Assets: no retained span contains a figure environment, a graphics-inclusion command, a PDF-page inclusion, a table or a TikZ picture; no image file is copied into this package, the package declares no asset dependency and no image file is redistributed with it. Licence: version arXiv:2405.17318v3 of this article is distributed under the Creative Commons Attribution 4.0 International licence, as recorded in the arXiv licence metadata for this exact version and shown on the abstract page https://arxiv.org/abs/2405.17318v3. A full rights notice is delivered at licenses/arxiv-2405.17318-CC-BY-4.0.txt. Characters: every retained excerpt is pure ASCII, so the delivered engine, XeLaTeX with its default Latin Modern text font, reports no missing character.
\begin{proposition} \label{p:sglim}
Let $[X,Y]^{\top}$ be a regularly varying random element in $L^2 \times L^2$. Then,
\[
\sg_{XY} = \lim_{n\to \infty}E\lb \lip \frac{X}{b(n)}, \frac{Y}{b(n)} \rip \Bigg|\|X\| \vee \|Y\| >b(n)\rb.
\]
\end{proposition}

\begin{assumption} \label{a:L2L2}
The bivariate random element $[X,Y]^{\top}$ in  $L^2 \times L^2$ has mean zero and is regularly varying with index $-\ag$, $\ag >2$. The observations $[X_1,Y_1]^{\top}, [X_2,Y_2]^{\top}, \ldots$ are independent copies of $[X,Y]^{\top}$.
\end{assumption}

\begin{lemma} \label{l:alM}
Under Assumption~\ref{a:L2L2}, for any $M>0$,
\[
\frac{n}{k} E \lb \lp \frac{R}{b(n/k)} \rp^2 I_{R \ge Mb(n/k) }  \rb  \to \frac{\ag}{\ag-2}M^{2-\ag}.
\]
\end{lemma}

\begin{lemma} \label{l:Bern}
Let $\bZ_n=(Z_1, \ldots, Z_n)$ with the $Z_i$ taking values in a
Lebesgue measurable subset $\cZ$ of an Euclidean space. Let $f$ be a real-valued function defined on $\cZ^n$. For $(z_1, \ldots, z_i) \in \cZ^{i}$, $1 \le i \le n$, put
\begin{equation} \label{eq:gi}
g_i(z_1, \ldots, z_i):=E\lb f(\bZ_n)| Z_j=z_j, 1 \le j \le i\rb -E\lb f(\bZ_n)| Z_j=z_j, 1 \le j \le i-1\rb.
\end{equation}
Define the maximum deviation by
\begin{equation} \label{eq:b}
b := \max_{1\le i\le n}\sup_{(z_1, \ldots, z_i) \in \cZ^i} g_i(z_1, \ldots, z_i),
\end{equation}
and define the supremum sum of variances by
\begin{equation} \label{eq:v}
\hat{v} := \sup_{(z_1, \ldots, z_n) \in \cZ^n} \sum_{i=1}^n \var\lb g_i(z_1, \ldots, z_{i-1}, Z_i^{\prime})\rb,
\end{equation}
where $Z_i^{\prime}$ is an independent copy of $Z_i$ conditional on $Z_j=z_j$, $1 \le j \le i-1$. If $b$ and $\hat{v}$ are finite, then for any $\eg \ge 0$,
\[
{\rm pr} \lp f(\bZ_n) - E[f(\bZ_n)] \ge t \rp \le \exp \lp \frac{-\eg^2}{2(\hat{v}+b\eg/3)}  \rp.
\]
\end{lemma}

\begin{proposition} \label{p:bar}
Under Assumption~\ref{a:L2L2},
\[
\sg_{n,k} \convP \sg_{XY}.
\]
\end{proposition}

\begin{proof}

Set
\begin{equation} \label{eq:sgbar}
\bar{\sg}_{n,k} := E\lb \lip \frac{X_1}{b(n/k)}, \frac{Y_1}{b(n/k)} \rip \Bigg|\|X_1\| \vee \|Y_1\| >b(n/k)\rb.
\end{equation}
Then, by Proposition~\ref{p:sglim}, $\bar{\sg}_{n,k} \to \sg_{XY}$, so it remains to show that $|\sg_{n,k} - \bar{\sg}_{n,k}| \convP 0$.

Let $\bZ_n = (Z_1, \ldots, Z_n)$, where $Z_i = (X_i, Y_i)$, and $\bz_n = (z_1, \ldots, z_n)$, where $z_i = (x_i, y_i)$, for $1 \le i \le n$. Consider a map $f: (L^2 \times L^2)^n \to \mbR$ defined by
\[
f(\bz_n) := \lmo \frac{1}{k}\sum_{i=1}^n \lip \frac{x_i}{b(n/k)}, \frac{y_i}{b(n/k)} \rip I_{r_i \ge b(n/k) }- \frac{n}{k}E\lb \lip \frac{X_1}{b(n/k)}, \frac{Y_1}{b(n/k)} \rip I_{R_1 >b(n/k)}\rb \rmo.
\]
Then, we have that
\[
|\sg_{n,k} - \bar{\sg}_{n,k}| = f(\bZ_n) - E[f(\bZ_n)] + E[f(\bZ_n)].
\]
We aim to show that $f(\bZ_n) - E[f(\bZ_n)] \convP 0$ and $E[f(\bZ_n)] \to 0$.

To establish the convergence, $f(\bZ_n) - E[f(\bZ_n)] \convP 0$, we use the Bernstein type concentration inequality in Lemma~\ref{l:Bern}. Since the $(X_i,Y_i)$ are independent, the deviation function in~\eqref{eq:gi} has the following form
\[
g_i(z_1, \ldots, z_i)= E \lb f(z_1, \ldots, z_{i-1},z_i, Z_{i+1}, \ldots, Z_n)-f(z_1, \ldots, z_{i-1},Z_i, Z_{i+1}, \ldots, Z_n) \rb.
\]
Then, using the fact that $\lmo |x|- |y| \rmo \le |x-y|$, we have that
\begin{align*}
g_i(z_1, \ldots, z_i)
&\le \frac{1}{k} E\lb \lmo  \lip \frac{x_i}{b(n/k)}, \frac{y_i}{b(n/k)} \rip I_{r_i \ge b(n/k) } - \lip \frac{X_i}{b(n/k)}, \frac{Y_i}{b(n/k)} \rip I_{R_i \ge b(n/k) }       \rmo \rb \\
&\le \frac{1}{k} \lbr \frac{|\lip x_i, y_i\rip|}{b(n/k)^2} + \frac{k}{n} \frac{n}{k} E \lb \lp \frac{R_i}{b(n/k)} \rp^2 I_{R_i \ge b(n/k) }  \rb \rbr\\
&\le \frac{1}{k} \lbr \frac{\|x_i\| \| y_i\|}{b(n/k)^2} + \frac{n}{k} E \lb \lp \frac{R_i}{b(n/k)} \rp^2 I_{R_i \ge b(n/k) }  \rb \rbr.
\end{align*}
Since $(x_i,y_i) \in L^2\times L^2$ and $\frac{n}{k} E \lb \lp {R_i}/{b(n/k)} \rp^2 I_{R_i \ge b(n/k) }  \rb \to \ag/(\ag-2)$ by Lemma~\ref{l:alM},  we have that $g_{i}(z_1, \ldots, z_i) \le {c_1}/{k}$, for some constant $c_1>0$. Therefore, the maximum deviation $b$ in~\eqref{eq:b} is bounded by $c_1/k$.

Next we investigate the upper bound for the sum of variances $\hat{v}$ in~\eqref{eq:v}. Since $E[g_i(z_1, \ldots, z_{i-1}, Z_i^{\prime})]=0$ by the law of total probability, we have that
\begin{align*}
&\var\lb g_i(z_1, \ldots, z_{i-1}, Z_i^{\prime})\rb \\
&= E [g_i^2(z_1, \ldots, z_{i-1}, Z_i^{\prime})]\\
&= E \lb \lbr  f(z_1, \ldots, z_{i-1},Z_i^{\prime}, Z_{i+1}, \ldots, Z_n)-f(z_1, \ldots, z_{i-1},Z_i, Z_{i+1}, \ldots, Z_n) \rbr^2 \rb \\
&\le \frac{1}{k^2} E\lb \lbr  \lip \frac{X_i^{\prime}}{b(n/k)}, \frac{Y_i^{\prime}}{b(n/k)} \rip I_{R_i^{\prime} \ge b(n/k) } - \lip \frac{X_i}{b(n/k)}, \frac{Y_i}{b(n/k)} \rip I_{R_i \ge b(n/k) }       \rbr^2 \rb \\
&\le \frac{2}{k^2} E \lb \lip \frac{X_i}{b(n/k)}, \frac{Y_i}{b(n/k)} \rip^2 I_{R_i \ge b(n/k) }   \rb \\
&\le \frac{2}{k^2} \lbr  \frac{k}{n} \frac{n}{k} E \lb  \lp \frac{R_i}{b(n/k)} \rp^2  I_{R_i \ge b(n/k) }  \rb \rbr.
\end{align*}
It then again follows from Lemma~\ref{l:alM} that $\var\lb g_i(z_1, \ldots, z_{i-1}, Z_i^{\prime})\rb  \le c_2/(nk)$ for some $c_2>0$. Then the supremum sum of variances $\hat{v}$ is bounded above by $c_2/k$. Therefore by Lemma~\ref{l:Bern}, for any $\eg>0$
\[
{\rm pr} \lp f(\bZ_n) - E[f(\bZ_n)]  \ge \eg \rp \le \exp \lp \frac{-k\eg^2}{c_1+c_2\eg/3} \rp.
\]
If we apply this inequality to $-f(\bZ_n)$, then we obtain the following `two-sided' inequality
\[
{\rm pr} \lp |f(\bZ_n) - E[f(\bZ_n)]|  \ge \eg \rp \le 2\exp \lp \frac{-k\eg^2}{c_1+c_2\eg/3} \rp.
\]
From this, we obtain that  $f(\bZ_n) - E[f(\bZ_n)] \convP 0$.

Next, to show $E[f(\bZ_n)] \to 0$, we set, for $1 \le i \le n$
\[
\Delta_i =  \lip \frac{X_i}{b(n/k)}, \frac{Y_i}{b(n/k)} \rip I_{R_i \ge b(n/k) }- E\lb \lip \frac{X_1}{b(n/k)}, \frac{Y_1}{b(n/k)} \rip I_{R_1 >b(n/k)}\rb.
\]
Then, we have that
\begin{align*}
E \lb f(\bZ_n) \rb = \frac{n}{k} E \lb \lmo \frac{1}{n} \sum_{i=1}^n \Delta_i \rmo \rb &\le \frac{n}{k} \lbr E \lb \lp \frac{1}{n} \sum_{i=1}^n \Delta_i \rp^2 \rb \rbr^{1/2}\\
&= \frac{n}{k} \lbr E \lb   \frac{1}{n^2} \sum_{i=1}^n \Delta_i^2 + \frac{1}{n^2} \sum_{i \neq j} \Delta_i\Delta_j \rb \rbr^{1/2}.
\end{align*}
Since the $\Delta_i$ are independent, $E[\Delta_i\Delta_j]=0$, for $i \neq j$. Therefore,
\begin{align*}
&E \lb f(\bZ_n) \rb \\
&\le \frac{n^{1/2}}{k} \lbr E \lb    \Delta_1^2  \rb \rbr^{1/2} \\
&= \frac{n^{1/2}}{k} \lbr  E  \lb \lp \lip \frac{X_1}{b(n/k)}, \frac{Y_1}{b(n/k)} \rip I_{R_1 \ge b(n/k) } - E\lb \lip \frac{X_1}{b(n/k)}, \frac{Y_1}{b(n/k)} \rip I_{R_1 >b(n/k)}\rb  \rp^2\rb \rbr^{1/2}\\
& = \frac{n^{1/2}}{k} \lbr  \var \lb  \lip \frac{X_1}{b(n/k)}, \frac{Y_1}{b(n/k)} \rip I_{R_1 \ge b(n/k) } \rb \rbr^{1/2} \\
& \le \frac{n^{1/2}}{k} \lbr  E  \lb \lip \frac{X_1}{b(n/k)}, \frac{Y_1}{b(n/k)} \rip^2 I_{R_1 \ge b(n/k) } \rb \rbr^{1/2}\\
& \le  \frac{n^{1/2}}{k}  \lbr E \lb  \lp \frac{R_1}{b(n/k)} \rp^2  I_{R_1 \ge b(n/k) }  \rb \rbr^{1/2}.
\end{align*}
Therefore, by Lemma~\ref{l:alM} we have that
\[
E \lb f(\bZ_n) \rb \le \frac{n^{1/2}}{k}  \lbr  \frac{k}{n} \frac{n}{k} E \lb  \lp \frac{R_1}{b(n/k)} \rp^2  I_{R_1 \ge b(n/k) }  \rb \rbr^{1/2} \le \frac{c_3}{k^{1/2}},
\]
for some $c_3>0$, which completes the proof.
\end{proof}

\end{document}
